% ------------------------------------------------------------------------------
% Capítulo · Conclusiones
% Fuentes: docs/tfg/{PROGRESO,BITACORA,TRAZABILIDAD,DECISIONES}.md.
% Estado de las fuentes: 20/10/2026 (F0–F3 cerradas; F4–F7 pendientes).
% Las filas marcadas como parciales o pendientes deben actualizarse al cerrar
% las fases F4–F7.
% ------------------------------------------------------------------------------

Este capítulo valora el grado de cumplimiento de los objetivos planteados en
la introducción, recoge las lecciones aprendidas durante el proyecto y
enumera las líneas de trabajo futuro.

\section{Objetivos alcanzados}
\label{sec:concl-objetivos}

El objetivo general era diseñar, implementar y validar una arquitectura
software reutilizable para el desarrollo ágil de aplicaciones SaaS orientadas
a pymes, con \textit{frontend}, \textit{backend}, autenticación, gestión de
usuarios y roles, pagos y un CMS para la administración de contenidos y datos.
PymeKit reúne todos esos componentes en un único monorepo y un único servicio
web, sobre una sola base de datos en la que se aplica la autorización. La
tabla~\ref{tab:concl-objetivos} detalla el resultado de cada objetivo
específico (tabla~\ref{tab:intro-objetivos}).

\todo{Actualizar el estado de OE-4, OE-7, OE-8 y OE-9 al cerrar las fases
F4--F7.}

\begin{table}[htbp]
  \centering
  \footnotesize
  \caption{Grado de cumplimiento de los objetivos específicos.}
  \label{tab:concl-objetivos}
  \begin{tabularx}{\textwidth}{@{}lX>{\raggedright\arraybackslash}p{2cm}@{}}
    \toprule
    \textbf{Obj.} & \textbf{Resultado} & \textbf{Estado} \\
    \midrule
    OE-1 & Contexto analizado, componentes comunes identificados y requisitos definidos: 13 funcionales, 9 no funcionales y 7 casos de uso, trazados al código, a las pruebas y a la memoria. & Cumplido \\
    OE-2 & Monorepo de paquetes desacoplados con exportaciones separadas de cliente y servidor; la API del CMS se monta dentro del servidor web, sin servicios adicionales (ADR-011). & Cumplido \\
    OE-3 & Registro, verificación de correo, contraseña, enlace mágico y MFA; cuentas personales y de equipo con invitaciones, roles y permisos; MFA obligatorio en la consola. Autorización aplicada en la base de datos mediante RLS. & Cumplido; roles adaptados a pymes pendientes (F5) \\
    OE-4 & Suscripciones con Stripe (planes, \textit{checkout}, portal del cliente y \textit{webhooks}) operativas en modo de pruebas. & Parcial: planes de ejemplo y pruebas E2E de pago pendientes (F5, F6) \\
    OE-5 & CMS de datos integrado en la consola: explorador de tablas con CRUD, relaciones y vistas guardadas, gestión de usuarios y archivos, auditoría no falsificable, búsqueda global, RBAC propio y paneles configurables (RF-08 a RF-11). & Cumplido \\
    OE-6 & Modelo multicliente con tablas compartidas y RLS; más de dos mil pruebas pgTAP, incluidas pruebas de ataque entre clientes; las auditorías posteriores a cada corrección concluyeron que el aislamiento se mantiene. & Cumplido \\
    OE-7 & Configuración centralizada en ficheros \texttt{*.config.ts} y variables \texttt{VITE\_*}; consola por áreas configurable desde la propia interfaz; interfaz solo en español con la traducción como punto de extensión documentado. & Parcial: documentación de la configuración (F5) y comentarios didácticos (F4) pendientes \\
    OE-8 & Pruebas unitarias, de base de datos y E2E en integración continua. & Parcial: E2E de CU-01 a CU-06 y evaluación de la reutilización (F6) pendientes \\
    OE-9 & Memoria y manuales redactados; configuración de despliegue en la nube preparada. & En curso (F7, F8) \\
    \bottomrule
  \end{tabularx}
\end{table}

% PENDIENTE (P-06): valoración de la aportación propia frente a la reutilizada.
Más allá de la tabla, el resultado principal del trabajo es haber convertido
dos bases de código independientes, con pilas, sesiones y modelos de permisos
distintos, en un único producto coherente: un solo inicio de sesión, una sola
consola, una sola forma de escribir formularios y textos, y una sola base de
datos cuya seguridad se ha auditado y reforzado. La integración del CMS en la
consola ---la decisión de mayor riesgo del proyecto--- se completó sin
recurrir al recorte previsto.

\section{Lecciones aprendidas}
\label{sec:concl-lecciones}

\subsection{Esfuerzo dedicado}

Hasta la fecha de corte (20/10/2026) el repositorio registra 54
\textit{commits}, 21 decisiones de arquitectura y 55 incidencias documentadas
en la bitácora. Las fases F0 a F3 ---arnés, base SaaS, integración completa
del CMS, desmarcado y español--- se ejecutaron en tres jornadas de trabajo con
un total de 23 horas del autor (tabla~\ref{tab:plan-dedicacion}). Al cierre de
la F2 la base de datos contaba con más de dos mil pruebas pgTAP y la suite E2E
superaba las 140 pruebas.
\todo{Completar con la dedicación total del proyecto y las cifras finales de
pruebas al cerrar F6.}

\subsection{Sobre reutilizar código de terceros}

\begin{itemize}
  \item \textbf{Reutilizar no exime de verificar.} De las 55 incidencias, 29
        son de seguridad y 25 de ellas venían del código de partida: escalada
        de privilegios por un \texttt{UPDATE} sin \texttt{WITH CHECK} (B-05),
        comprobaciones de MFA y de bloqueo que fallaban en abierto (B-06,
        B-23), lectura de tablas internas (B-07, B-08, B-28), inyección SQL
        (B-24), auditoría falsificable (B-27) o escape de denegaciones por
        delegación (B-47). Una base reutilizable solo acelera si se ha
        verificado antes; en caso contrario, multiplica sus fallos en cada
        proyecto que la usa.
  \item \textbf{Un test en verde no demuestra corrección.} Un test heredado
        afirmaba precisamente el comportamiento inseguro (B-11). Hay que leer
        las aserciones, no solo los nombres de las pruebas.
  \item \textbf{El inventario de lo que se porta debe hacerse ruta a ruta.} La
        búsqueda global del CMS no figuraba en el plan inicial porque el
        inventario se hizo por paquetes (B-30).
  \item \textbf{El desmarcado es más que un \texttt{grep}.} La marca aparecía
        en expresiones regulares escapadas (B-01), en nombres de carpetas y
        rutas de importación (B-40) y en imágenes que ningún control textual
        puede leer (B-46). Los controles automáticos deben cubrir todas las
        variantes, y aun así hace falta una revisión visual.
\end{itemize}

\subsection{Sobre seguridad en la base de datos}

\begin{itemize}
  \item \textbf{Fallar en cerrado.} Toda comprobación de seguridad que
        encuentra un error debe denegar (B-06, B-23).
  \item \textbf{«Sin error» no es «éxito».} Una escritura filtrada por RLS
        afecta a cero filas sin producir ningún error, y hay que comprobarlo
        (B-49).
  \item \textbf{Defensa en profundidad.} Las reglas deben aplicarse en la
        base de datos y no solo en la API (B-31); una operación pensada para un
        registro debe comprobar que su condición identifica exactamente uno
        (B-19).
  \item \textbf{Las herramientas de diferencias no ven la seguridad.} La
        generación automática de migraciones omite opciones como
        \texttt{security\_invoker} de las vistas y los privilegios por columna
        (B-35); esas propiedades se verifican con pruebas, no con la
        herramienta.
  \item \textbf{«Local» no significa privado.} En un servidor con IP pública,
        el motor de contenedores puede saltarse el cortafuegos (B-04).
\end{itemize}

\subsection{Sobre el desarrollo asistido por IA}

\begin{itemize}
  \item \textbf{El cuello de botella se desplaza a la verificación.} La
        integración del CMS se estimó en semanas de desarrollo convencional y
        se ejecutó en días (sección~\ref{sec:plan-estimado-real}), pero buena
        parte del tiempo del autor se dedicó a revisar, probar y decidir.
  \item \textbf{La revisión humana sigue siendo necesaria.} El autor detectó
        con una prueba manual que la consola se cargaba sin segundo factor
        (B-14); una inyección SQL heredada pasó desapercibida en varias
        pasadas de los agentes por el mismo fichero (B-24), y otra latente,
        introducida al portar el código, solo la encontró una revisión
        posterior que buscaba otra cosa (B-45).
  \item \textbf{El nivel de revisión se decide por riesgo.} La verificación
        exhaustiva con varios agentes encontró brechas reales, pero su coste
        obligó a ajustarla (B-48); al reducirla, aparecieron pruebas
        desactualizadas en zonas no revisadas (B-51, B-54). La regla resultante
        es ejecutar siempre la suite completa antes de fusionar.
  \item \textbf{Las directrices escritas son el mecanismo de control.} Las
        reglas del arnés (leer el progreso antes de empezar, no reabrir
        decisiones sin un ADR, verificar antes de cerrar, registrar cada
        incidencia) son las que hacen que el trabajo de los agentes sea
        evaluable y que esta memoria pueda apoyarse en evidencias.
  \item \textbf{Cada incremento, su \textit{commit}.} Un corte de conexión
        con un agente trabajando dejó cambios de distinta naturaleza mezclados
        en el árbol de trabajo (B-15).
\end{itemize}

\todo{El autor añade su valoración personal: qué ha aprendido, qué haría de
otra manera y cómo ha cambiado su forma de trabajar con agentes de IA.}

\section{Trabajo futuro}
\label{sec:concl-futuro}

Las líneas de trabajo futuro se dividen entre lo que queda del plan del
proyecto y las mejoras detectadas durante su desarrollo.

\paragraph{Fases pendientes del plan.}
\begin{itemize}
  \item \textbf{F4, comentarios didácticos}: reescribir en español los
        comentarios heredados de todos los módulos, empezando por los
        esquemas de la base de datos y el \textit{middleware} de
        autorización.
  \item \textbf{F5, adaptación a pymes}: roles y permisos de equipo
        orientados a pymes, planes de suscripción de ejemplo, una plantilla de
        permisos inicial del CMS (soporte y gestor de contenidos) y la
        documentación completa de la configuración reutilizable.
  \item \textbf{F6, validación}: pruebas E2E de los casos de uso CU-01 a
        CU-06, pruebas de integración con los \textit{webhooks} de Stripe y la
        evaluación de la reutilización (CU-07).
  \item \textbf{F7, despliegue}: despliegue de demostración en la nube y
        manual de instalación probado desde cero.
\end{itemize}

\paragraph{Mejoras del producto.}
\begin{itemize}
  \item Completar el CMS: renombrar carpetas en el explorador de archivos,
        editores de texto enriquecido y de ficheros en los formularios, una
        interfaz para definir relaciones virtuales y ajustes por tipo de
        formato (moneda, fichero).
  \item Cerrar los pendientes heredados documentados: el total incorrecto de
        la lista de paneles y la denegación de almacenamiento sobre una
        subruta.
  \item Plantillas de correo en español para todos los mensajes y fechas
        del CMS según el idioma activo.
  \item Estabilizar las pruebas E2E inestables bajo carga.
  \item Ofrecer otros idiomas siguiendo el punto de extensión documentado.
  \item Un generador que cree un proyecto nuevo a partir de PymeKit con su
        configuración inicial, para reducir aún más el tiempo de puesta en
        marcha.
  \item Explorar una variante de PymeKit pensada para dar servicio a muchas
        pymes desde una sola instalación, apoyada en las cuentas de equipo
        que ya ofrece.
\end{itemize}
